%%%%%%%%%%%%%%%%%%%%%%%%%%%%%%%%%%%%%%%%%%%%%%%%%%%%%%%%%%%%%
%% Begin exercise %%
%%%%%%%%%%%%%%%%%%%%%%%%%%%%%%%%%%%%%%%%%%%%%%%%%%%%%%%%%%%%%
\ex{Frequency scaling of power converters}

%%%%%%%%%%%%%%%%%%%%%%%%%%%%%%%%%%%%%%%%%%%%%%%%%%%%%%%%%%%%%
%% Task 1 %%
%%%%%%%%%%%%%%%%%%%%%%%%%%%%%%%%%%%%%%%%%%%%%%%%%%%%%%%%%%%%%

\task{Passive components of a buck converter versus switching frequency}

The synchronous buck converter in \autoref{fig:ex01_buck} operates in continuous conduction mode with the data in \autoref{table:ex01_buck}. All components are ideal.

%%%%%%%%%%%%%%%%%%%%%%%%%%%%%%%%%%%%%%%%%%%%%%%%%%%%%%%%%%%%%%%%%%%%%%%%%%
% Synchronous buck converter (Exercise 01, Task 1)
%%%%%%%%%%%%%%%%%%%%%%%%%%%%%%%%%%%%%%%%%%%%%%%%%%%%%%%%%%%%%%%%%%%%%%%%%%

\begin{figure}[htb]
	\begin{center}
		\begin{circuitikz}
			% transistors
			\draw (2.5,2.6) node[nigfete,rotate=90](S1){} (3.8,1.2) node[nigfete](S2){};
			\draw (S1) node[anchor=south, yshift=0.3cm]{$S_1$};
			\draw (S2) node[anchor=west]{$S_2$};
			% input source and high-side switch
			\draw (0,0) to[V, l=$V_\mathrm{in}$] (0,2.6) -- (S1.D);
			\draw (S1.S) -- (3.8,2.6) coordinate (sw);
			% low-side switch from switch node to ground
			\draw (sw) to[short, *-] (S2.D);
			\draw (S2.S) -- (3.8,0);
			% inductor, capacitor, load
			\draw (sw) to[L, l=$L$, i>^=$i_L$] (6.3,2.6) to[short, -*] (6.8,2.6) coordinate (out) -- (8.4,2.6) to[R, l=$R$] (8.4,0) -- (0,0);
			\draw (out) to[C, l_=$C$] (6.8,0) node[circ]{};
			\draw (3.8,0) node[circ]{};
			\draw (9.3,2.6) to[open, v^=$V_\mathrm{o}$] (9.3,0);
		\end{circuitikz}
	\end{center}
	\caption{Synchronous buck converter.}
	\label{fig:ex01_buck}
\end{figure}


\begin{table}[ht]
	\centering
	\begin{tabular}{ll}
		\toprule
		Input voltage $V_\mathrm{in}$                       & $\SI{24}{\volt}$                       \\
		Output voltage $V_\mathrm{o}$                       & $\SI{5}{\volt}$                        \\
		Output current $I_\mathrm{o}$                       & $\SI{6}{\ampere}$                      \\
		Peak-to-peak inductor current ripple $\Delta I_L$   & $\SI{40}{\percent}$ of $I_\mathrm{o}$  \\
		Peak-to-peak output voltage ripple $\Delta V_\mathrm{o}$ & $\SI{25}{\milli\volt}$            \\
		\bottomrule
	\end{tabular}
	\caption{Specification of the buck converter.}
	\label{table:ex01_buck}
\end{table}

\subtask{Calculate the duty cycle $D$, the inductance $L$ and the capacitance $C$ for $f_\mathrm{sw}=\SI{200}{\kilo\hertz}$, $\SI{2}{\mega\hertz}$ and $\SI{20}{\mega\hertz}$.}

\begin{solutionblock}
	The duty cycle in continuous conduction mode is
	\begin{equation}
		D = \frac{V_\mathrm{o}}{V_\mathrm{in}} = \frac{\SI{5}{\volt}}{\SI{24}{\volt}} = 0.208.
	\end{equation}
	With $\Delta I_L = 0.4\cdot\SI{6}{\ampere} = \SI{2.4}{\ampere}$ the required inductance and capacitance are
	\begin{equation}
		L = \frac{V_\mathrm{o}(1-D)}{f_\mathrm{sw}\,\Delta I_L}, \qquad C = \frac{\Delta I_L}{8 f_\mathrm{sw}\,\Delta V_\mathrm{o}}.
	\end{equation}
	The results are listed in \autoref{stable:ex01_LC}. Both $L$ and $C$ scale with $1/f_\mathrm{sw}$.
	\begin{solutiontable}[ht]
		\centering
		\begin{tabular}{lccc}
			\toprule
			$f_\mathrm{sw}$ & $\SI{200}{\kilo\hertz}$  & $\SI{2}{\mega\hertz}$   & $\SI{20}{\mega\hertz}$  \\
			\midrule
			$L$             & $\SI{8.25}{\micro\henry}$ & $\SI{825}{\nano\henry}$ & $\SI{82.5}{\nano\henry}$ \\
			$C$             & $\SI{60}{\micro\farad}$   & $\SI{6}{\micro\farad}$  & $\SI{0.6}{\micro\farad}$ \\
			\bottomrule
		\end{tabular}
		\caption{Inductance and capacitance versus switching frequency.}
		\label{stable:ex01_LC}
	\end{solutiontable}
\end{solutionblock}

\subtask{Calculate the peak energy stored in the inductor and compare it with the energy $W=P_\mathrm{o}/f_\mathrm{sw}$ delivered to the load in one switching period. What do you conclude?}

\begin{solutionblock}
	The peak inductor current is $\hat{i}_L = I_\mathrm{o}+\Delta I_L/2 = \SI{7.2}{\ampere}$ and the output power is $P_\mathrm{o}=\SI{30}{\watt}$.
	\begin{equation}
		E_L = \tfrac{1}{2} L \hat{i}_L^2 = \SI{214}{\micro\joule},\ \SI{21.4}{\micro\joule},\ \SI{2.14}{\micro\joule}
		\qquad
		W = \frac{P_\mathrm{o}}{f_\mathrm{sw}} = \SI{150}{\micro\joule},\ \SI{15}{\micro\joule},\ \SI{1.5}{\micro\joule}
	\end{equation}
	for $\SI{200}{\kilo\hertz}$, $\SI{2}{\mega\hertz}$ and $\SI{20}{\mega\hertz}$. The ratio
	\begin{equation}
		\frac{E_L}{W} = \frac{(1-D)\,\hat{i}_L^2}{2\,\Delta I_L\,I_\mathrm{o}} = 1.43
	\end{equation}
	does not depend on $f_\mathrm{sw}$. The inductor always stores about 1.4 times the energy processed per period, and this energy falls with $1/f_\mathrm{sw}$. A smaller stored energy allows a smaller magnetic component only if the flux density and the current density can be kept, which is the subject of the next subtask.
\end{solutionblock}

\subtask{Give three reasons why the inductor at $\SI{20}{\mega\hertz}$ is not 100 times smaller than the one at $\SI{200}{\kilo\hertz}$.}

\begin{solutionblock}
	\begin{itemize}
		\item The winding still carries $\SI{6}{\ampere}$ DC. The conductor cross-section is set by the current, not by $f_\mathrm{sw}$.
		\item Core loss density rises with $f^{\alpha}$ ($\alpha\approx1.4$--$2$ for ferrites). The peak AC flux density must therefore fall with frequency, so the core cannot shrink in proportion to $L$.
		\item The skin depth falls from $\SI{148}{\micro\meter}$ at $\SI{200}{\kilo\hertz}$ to $\SI{15}{\micro\meter}$ at $\SI{20}{\mega\hertz}$. The AC resistance of the winding rises, and proximity effect in multilayer windings makes it worse.
		\item Terminations, clearances and the minimum practical size of cores and bobbins do not scale.
	\end{itemize}
\end{solutionblock}

%%%%%%%%%%%%%%%%%%%%%%%%%%%%%%%%%%%%%%%%%%%%%%%%%%%%%%%%%%%%%
%% Task 2 %%
%%%%%%%%%%%%%%%%%%%%%%%%%%%%%%%%%%%%%%%%%%%%%%%%%%%%%%%%%%%%%

\task{Corner frequency of an input EMI filter}

A converter needs an input filter that attenuates the fundamental switching component by $A=\SI{60}{\decibel}$. Use an ideal second-order LC low-pass filter whose attenuation rises with $\SI{40}{\decibel}$ per decade above its corner frequency $f_\mathrm{c}$.

\subtask{Derive $f_\mathrm{c}$ as a function of $f_\mathrm{sw}$ and $A$. Calculate $f_\mathrm{c}$ and the product $LC$ for $f_\mathrm{sw}=\SI{200}{\kilo\hertz}$ and $f_\mathrm{sw}=\SI{2}{\mega\hertz}$.}

\begin{solutionblock}
	From $A = 40\log_{10}\left(f_\mathrm{sw}/f_\mathrm{c}\right)$ follows
	\begin{equation}
		f_\mathrm{c} = \frac{f_\mathrm{sw}}{10^{A/40}} = \frac{f_\mathrm{sw}}{31.6}, \qquad LC = \frac{1}{\left(2\pi f_\mathrm{c}\right)^2}.
	\end{equation}
	For $\SI{200}{\kilo\hertz}$: $f_\mathrm{c}=\SI{6.32}{\kilo\hertz}$ and $LC=\SI{6.33 \cdot 10^{-10}}{\square\second}$, e.g.\ $\SI{10}{\micro\henry}$ with $\SI{63.3}{\micro\farad}$.\\
	For $\SI{2}{\mega\hertz}$: $f_\mathrm{c}=\SI{63.2}{\kilo\hertz}$ and $LC=\SI{6.33 \cdot 10^{-12}}{\square\second}$, e.g.\ $\SI{1}{\micro\henry}$ with $\SI{6.33}{\micro\farad}$.\\
	Ten times the switching frequency allows a 100 times smaller $LC$ product.
\end{solutionblock}

\subtask{Many flyback adapters switch at about $\SI{65}{\kilo\hertz}$. Explain this choice with respect to the conducted-emission band from $\SI{150}{\kilo\hertz}$ to $\SI{30}{\mega\hertz}$.}

\begin{solutionblock}
	The fundamental at $\SI{65}{\kilo\hertz}$ and the second harmonic at $\SI{130}{\kilo\hertz}$ lie below $\SI{150}{\kilo\hertz}$ and are not regulated. The first harmonic inside the band is the third one at $\SI{195}{\kilo\hertz}$. It is smaller than the fundamental, and the filter has to provide its attenuation at a higher frequency. Both relax the filter.
\end{solutionblock}

\subtask{Why does the filter at $\SI{2}{\mega\hertz}$ not become 100 times smaller in practice?}

\begin{solutionblock}
	The ideal $-\SI{40}{\decibel}$ per decade holds only up to the self-resonance of the filter components. The winding capacitance of the inductor and the ESL of the capacitor limit the attenuation at high frequency, and the regulated band extends to $\SI{30}{\mega\hertz}$. In addition, the components must still carry the input current and withstand the input voltage, and common-mode noise from $C_\mathrm{p}\,\mathrm{d}v/\mathrm{d}t$ needs its own filter stage.
\end{solutionblock}

%%%%%%%%%%%%%%%%%%%%%%%%%%%%%%%%%%%%%%%%%%%%%%%%%%%%%%%%%%%%%
%% Task 3 %%
%%%%%%%%%%%%%%%%%%%%%%%%%%%%%%%%%%%%%%%%%%%%%%%%%%%%%%%%%%%%%

\task{Skin depth and winding design}

Use the resistivity of copper $\rho_{20}=\SI{1.72 \cdot 10^{-8}}{\ohm\meter}$ at $\SI{20}{\celsius}$, its temperature coefficient $\alpha_\mathrm{Cu}=\SI{3.93 \cdot 10^{-3}}{\per\kelvin}$ and $\mu_0=4\pi\cdot\SI{10^{-7}}{\henry\per\meter}$.

\subtask{Calculate the skin depth $\delta=\sqrt{\rho/(\pi f \mu_0)}$ of copper at $\SI{20}{\celsius}$ and at $\SI{100}{\celsius}$ for $\SI{100}{\kilo\hertz}$, $\SI{1}{\mega\hertz}$, $\SI{6.78}{\mega\hertz}$, $\SI{13.56}{\mega\hertz}$ and $\SI{30}{\mega\hertz}$.}

\begin{solutionblock}
	At $\SI{100}{\celsius}$ the resistivity is $\rho_{100}=\rho_{20}\,(1+\alpha_\mathrm{Cu}\cdot\SI{80}{\kelvin})=1.31\,\rho_{20}$, so $\delta$ is $\sqrt{1.31}=1.15$ times larger. The results are listed in \autoref{stable:ex01_skin}.
	\begin{solutiontable}[ht]
		\centering
		\begin{tabular}{lccccc}
			\toprule
			$f$                               & $\SI{100}{\kilo\hertz}$ & $\SI{1}{\mega\hertz}$ & $\SI{6.78}{\mega\hertz}$ & $\SI{13.56}{\mega\hertz}$ & $\SI{30}{\mega\hertz}$ \\
			\midrule
			$\delta$ at $\SI{20}{\celsius}$   & $\SI{209}{\micro\meter}$ & $\SI{66.0}{\micro\meter}$ & $\SI{25.3}{\micro\meter}$ & $\SI{17.9}{\micro\meter}$ & $\SI{12.1}{\micro\meter}$ \\
			$\delta$ at $\SI{100}{\celsius}$  & $\SI{239}{\micro\meter}$ & $\SI{75.7}{\micro\meter}$ & $\SI{29.1}{\micro\meter}$ & $\SI{20.6}{\micro\meter}$ & $\SI{13.8}{\micro\meter}$ \\
			\bottomrule
		\end{tabular}
		\caption{Skin depth of copper.}
		\label{stable:ex01_skin}
	\end{solutiontable}
\end{solutionblock}

\subtask{A round copper wire with a diameter of $\SI{1}{\milli\meter}$ carries a $\SI{1}{\mega\hertz}$ current. Estimate $R_\mathrm{ac}/R_\mathrm{dc}$ by assuming that the current flows uniformly in a layer of thickness $\delta$ below the surface. Compare with the exact value of 4.05.}

\begin{solutionblock}
	With the radius $r=\SI{0.5}{\milli\meter}$ and $\delta\ll r$, the conducting area is about $2\pi r\delta$ instead of $\pi r^2$:
	\begin{equation}
		\frac{R_\mathrm{ac}}{R_\mathrm{dc}} \approx \frac{\pi r^2}{2\pi r\delta} = \frac{r}{2\delta} = \frac{\SI{500}{\micro\meter}}{2\cdot\SI{66.0}{\micro\meter}} = 3.79.
	\end{equation}
	This is $\SI{6}{\percent}$ below the exact value from the Bessel-function solution. The asymptotic expansion of the exact solution adds $1/4$, giving $r/(2\delta)+1/4=4.04$.
\end{solutionblock}

\subtask{A single-layer foil winding carries a $\SI{6.78}{\mega\hertz}$ current. The magnetic field is on one side of the foil only. Its AC resistance per unit width and length is
	\begin{equation*}
		R'_\mathrm{ac} = \frac{\rho}{\delta}\,g(\Delta), \qquad g(\Delta)=\frac{\sinh 2\Delta + \sin 2\Delta}{\cosh 2\Delta - \cos 2\Delta}, \qquad \Delta = \frac{h}{\delta}
	\end{equation*}
	with the foil thickness $h$. Show that $R'_\mathrm{ac}$ is minimal for $h=\frac{\pi}{2}\delta$. Compare $h=\SI{35}{\micro\meter}$ (1\,oz) with $h=\SI{70}{\micro\meter}$ (2\,oz) at $\SI{20}{\celsius}$.}

\begin{solutionblock}
	With $x=2\Delta$ and the identities $\cosh^2 x-\sinh^2 x=1$ and $\cos^2 x+\sin^2 x=1$ the derivative becomes
	\begin{equation}
		\frac{\mathrm{d}g}{\mathrm{d}x} = \frac{(\cosh x+\cos x)(\cosh x-\cos x)-(\sinh x+\sin x)^2}{(\cosh x-\cos x)^2} = \frac{-2\sinh x\,\sin x}{(\cosh x-\cos x)^2}.
	\end{equation}
	It is negative for $0<x<\pi$ and positive for $\pi<x<2\pi$, so the minimum is at $x=\pi$, i.e.\ $h=\frac{\pi}{2}\delta$ with $g=0.917$. Thicker foil does not help: $g\rightarrow 1$ for $h\gg\delta$.\\
	At $\SI{6.78}{\mega\hertz}$, $\delta=\SI{25.3}{\micro\meter}$ and the optimum is $h=\SI{39.8}{\micro\meter}$.
	\begin{itemize}
		\item 1\,oz: $\Delta=1.38$, $g=0.930$, only $\SI{1.4}{\percent}$ above the optimum.
		\item 2\,oz: $\Delta=2.76$, $g=1.000$, i.e.\ $\SI{7.5}{\percent}$ more AC resistance than 1\,oz.
	\end{itemize}
	The thinner copper is the better choice. The extra copper of the 2\,oz foil carries almost no current.
\end{solutionblock}

\subtask{By how much does the AC resistance of the 1\,oz foil rise from $\SI{20}{\celsius}$ to $\SI{100}{\celsius}$? Compare with the DC resistance.}

\begin{solutionblock}
	The DC resistance rises with $\rho$, i.e.\ by $\SI{31}{\percent}$. The AC resistance is proportional to $\rho/\delta\cdot g(\Delta)\propto\sqrt{\rho}\,g(\Delta)$. At $\SI{100}{\celsius}$, $\delta=\SI{29.1}{\micro\meter}$, $\Delta=1.20$ and $g=0.974$:
	\begin{equation}
		\frac{R'_\mathrm{ac,100}}{R'_\mathrm{ac,20}} = \sqrt{1.31}\cdot\frac{0.974}{0.930} = 1.20.
	\end{equation}
	The AC resistance rises by only $\SI{20}{\percent}$, because the larger skin depth uses more of the copper.
\end{solutionblock}

%%%%%%%%%%%%%%%%%%%%%%%%%%%%%%%%%%%%%%%%%%%%%%%%%%%%%%%%%%%%%
%% Task 4 %%
%%%%%%%%%%%%%%%%%%%%%%%%%%%%%%%%%%%%%%%%%%%%%%%%%%%%%%%%%%%%%

\task{Capacitive switching loss and the limit of hard switching}

In the half-bridge of \autoref{fig:ex01_halfbridge}, $S_1$ turns on while it blocks the full voltage $V_\mathrm{DC}$ (hard turn-on). Both switches are identical. Their output charge $Q_\mathrm{oss}(V_\mathrm{DC})=\int_0^{V_\mathrm{DC}}C_\mathrm{oss}(v)\,\mathrm{d}v$ is given in \autoref{table:ex01_cases}. The values are illustrative and do not describe a specific device.

%%%%%%%%%%%%%%%%%%%%%%%%%%%%%%%%%%%%%%%%%%%%%%%%%%%%%%%%%%%%%%%%%%%%%%%%%%
% Half-bridge with output capacitances (Exercise 01, Task 4)
%%%%%%%%%%%%%%%%%%%%%%%%%%%%%%%%%%%%%%%%%%%%%%%%%%%%%%%%%%%%%%%%%%%%%%%%%%

\begin{figure}[htb]
	\begin{center}
		\begin{circuitikz}
			% transistors
			\draw (3,3.75) node[nigfete](S1){} (3,1.25) node[nigfete](S2){};
			\draw (S1) node[anchor=east, xshift=-0.9cm]{$S_1$};
			\draw (S2) node[anchor=east, xshift=-0.9cm]{$S_2$};
			% DC link
			\draw (0,0) to[V, l=$V_\mathrm{DC}$] (0,5) -- (3,5) -- (S1.D);
			\draw (S1.S) -- (S2.D);
			\draw (S2.S) -- (3,0) -- (0,0);
			% output capacitances in parallel to the switches
			\draw (3,5) -- (4.3,5) to[C, l=$C_\mathrm{oss,1}$] (4.3,2.5);
			\draw (3,2.5) to[short, *-*] (4.3,2.5) to[C, l=$C_\mathrm{oss,2}$] (4.3,0) -- (3,0);
			% load current
			\draw (4.3,2.5) to[short, i=$I_\mathrm{L}$, -o] (6.0,2.5);
			\draw (3,5) node[circ]{} (3,0) node[circ]{};
		\end{circuitikz}
	\end{center}
	\caption{Half-bridge with the output capacitances of both switches. Before $S_1$ turns on, the load current flows through $S_2$ in reverse direction.}
	\label{fig:ex01_halfbridge}
\end{figure}


\begin{table}[ht]
	\centering
	\begin{tabular}{llll}
		\toprule
		Application                         & $V_\mathrm{DC}$          & $Q_\mathrm{oss}(V_\mathrm{DC})$ & $P_\mathrm{o}$          \\
		\midrule
		A: $\SI{650}{\volt}$ GaN, PFC or DC-DC  & $\SI{400}{\volt}$    & $\SI{56}{\nano\coulomb}$        & $\SI{1.5}{\kilo\watt}$  \\
		B: $\SI{100}{\volt}$ GaN, bus converter & $\SI{48}{\volt}$     & $\SI{25}{\nano\coulomb}$        & $\SI{500}{\watt}$       \\
		C: $\SI{30}{\volt}$ GaN, point of load  & $\SI{12}{\volt}$     & $\SI{5}{\nano\coulomb}$         & $\SI{50}{\watt}$        \\
		\bottomrule
	\end{tabular}
	\caption{Three applications with hard-switched half-bridges.}
	\label{table:ex01_cases}
\end{table}

\subtask{Show that the energy dissipated at each hard turn-on is $E=Q_\mathrm{oss}V_\mathrm{DC}$, independent of the shape of $C_\mathrm{oss}(v)$. Neglect the $V$--$I$ overlap.}

\begin{solutionblock}
	Before the turn-on, $C_\mathrm{oss,1}$ is charged to $V_\mathrm{DC}$ and stores $E_\mathrm{oss}(V_\mathrm{DC})$, while $C_\mathrm{oss,2}$ is discharged.
	\begin{itemize}
		\item The channel of $S_1$ discharges $C_\mathrm{oss,1}$. Its stored energy $E_\mathrm{oss}$ is dissipated in $S_1$.
		\item The channel of $S_1$ also carries the current that charges $C_\mathrm{oss,2}$ to $V_\mathrm{DC}$. The DC link delivers the charge $Q_\mathrm{oss}$ and the energy $Q_\mathrm{oss}V_\mathrm{DC}$. Of this, $E_\mathrm{oss}$ remains stored in $C_\mathrm{oss,2}$, and $Q_\mathrm{oss}V_\mathrm{DC}-E_\mathrm{oss}$ is dissipated in $S_1$.
	\end{itemize}
	In total
	\begin{equation}
		E = E_\mathrm{oss} + \left(Q_\mathrm{oss}V_\mathrm{DC}-E_\mathrm{oss}\right) = Q_\mathrm{oss}V_\mathrm{DC}.
	\end{equation}
	Only the charge $Q_\mathrm{oss}(V_\mathrm{DC})$ appears, not the shape of $C_\mathrm{oss}(v)$.
\end{solutionblock}

\subtask{Calculate $E$ and the capacitive switching loss $P=E f_\mathrm{sw}$ at $\SI{100}{\kilo\hertz}$ and $\SI{1}{\mega\hertz}$ for applications A, B and C. At which $f_\mathrm{sw}$ does $P$ reach $\SI{1}{\percent}$ of $P_\mathrm{o}$?}

\begin{solutionblock}
	The results are listed in \autoref{stable:ex01_qoss}, with $f_{1\%}=0.01\,P_\mathrm{o}/(Q_\mathrm{oss}V_\mathrm{DC})$.
	\begin{solutiontable}[ht]
		\centering
		\begin{tabular}{lcccc}
			\toprule
			Application & $E=Q_\mathrm{oss}V_\mathrm{DC}$ & $P$ at $\SI{100}{\kilo\hertz}$ & $P$ at $\SI{1}{\mega\hertz}$ & $f_{1\%}$ \\
			\midrule
			A & $\SI{22.4}{\micro\joule}$ & $\SI{2.24}{\watt}$  & $\SI{22.4}{\watt}$ & $\SI{670}{\kilo\hertz}$ \\
			B & $\SI{1.2}{\micro\joule}$  & $\SI{0.12}{\watt}$  & $\SI{1.2}{\watt}$  & $\SI{4.2}{\mega\hertz}$ \\
			C & $\SI{60}{\nano\joule}$    & $\SI{6}{\milli\watt}$ & $\SI{60}{\milli\watt}$ & $\SI{8.3}{\mega\hertz}$ \\
			\bottomrule
		\end{tabular}
		\caption{Capacitive switching loss of the three applications.}
		\label{stable:ex01_qoss}
	\end{solutiontable}
	At $\SI{400}{\volt}$ hard switching ends below $\SI{1}{\mega\hertz}$, while low-voltage converters can be hard-switched into the MHz range. Above these frequencies the capacitive energy must be recovered by zero-voltage switching (L8--L10).
\end{solutionblock}

\subtask{The datasheet of the device in application A gives $C_\mathrm{oss}=\SI{50}{\pico\farad}$ at $\SI{400}{\volt}$. Calculate $\frac{1}{2}C_\mathrm{oss}V_\mathrm{DC}^2$ and compare it with $Q_\mathrm{oss}V_\mathrm{DC}$. Explain the difference.}

\begin{solutionblock}
	\begin{equation}
		\tfrac{1}{2}C_\mathrm{oss}V_\mathrm{DC}^2 = \tfrac{1}{2}\cdot\SI{50}{\pico\farad}\cdot(\SI{400}{\volt})^2 = \SI{4}{\micro\joule} \ll Q_\mathrm{oss}V_\mathrm{DC} = \SI{22.4}{\micro\joule}.
	\end{equation}
	$C_\mathrm{oss}$ is strongly nonlinear and much larger at low voltage, so most of $Q_\mathrm{oss}$ is stored at low voltage. The small-signal value at $\SI{400}{\volt}$ therefore underestimates the charge. Moreover, $\frac{1}{2}CV^2$ is a stored energy, while the loss of a hard turn-on is $Q_\mathrm{oss}V_\mathrm{DC}$. Use $Q_\mathrm{oss}$ or the charge-equivalent capacitance $C_\mathrm{o(tr)}=Q_\mathrm{oss}/V_\mathrm{DC}$ for this calculation (L2).
\end{solutionblock}

\vspace{0.5cm}
Remark: The notebook \texttt{tutorialEx01\_task.ipynb} extends Tasks 1.1 to 1.4 to frequency sweeps.
